\begin{table}[H]
\centering
\caption{Exploratory Credibility Moderation Relative to Apolitical Content}
\label{tab:credibility_moderation_t1}
\footnotesize
\begin{threeparttable}
\begin{tabular}{lcccc}
\toprule
 & Regime support & Censorship support & Nationalism & Trust in foreign media \\
\midrule
Pooled Pro-China & 0.169*** & 0.013 & 0.067* & 0.072* \\
 & (0.024) & (0.028) & (0.027) & (0.029) \\
\addlinespace
Pooled Anti-China & -0.305*** & -0.103*** & 0.119*** & -0.037 \\
 & (0.023) & (0.028) & (0.027) & (0.029) \\
\addlinespace
Assigned-content credibility (z) & 0.014 & -0.032 & -0.050* & 0.085** \\
 & (0.020) & (0.024) & (0.023) & (0.026) \\
\addlinespace
Pro x credibility & 0.039 & 0.039 & 0.037 & 0.073* \\
 & (0.024) & (0.029) & (0.027) & (0.032) \\
\addlinespace
Anti x credibility & -0.104*** & 0.017 & -0.011 & 0.053+ \\
 & (0.024) & (0.029) & (0.027) & (0.032) \\
\addlinespace
\midrule
Control mean & 0.195 & -0.002 & -0.009 & 2.825 \\
Observations & 4,348 & 4,348 & 4,348 & 4,348 \\
Lagged dependent variable & Yes & Yes & Yes & Yes \\
Block fixed effects & Yes & Yes & Yes & Yes \\
HC2 standard errors & Yes & Yes & Yes & Yes \\
\bottomrule
\end{tabular}
\begin{tablenotes}[flushleft]
\footnotesize
\item Note: Each column reports an endline ANCOVA specification estimated on the pooled Pro-China, pooled Anti-China, and apolitical-China groups only, so that the apolitical-China group is the omitted reference category. `Assigned-content credibility' is the participant's mean slot-1 credibility rating over assigned weekly materials, standardized within the non-control endline sample. Because credibility is post-treatment, these interactions should be interpreted as exploratory robustness checks rather than purely causal moderators.
\end{tablenotes}
\end{threeparttable}
\end{table}
